\chapter{Razredi in družine}
\label{cha:razred-druzine}

V tem poglavju spoznamo razrede, ki jih uvedemo, da razrešimo Russellov paradoks. Nato spoznamo še družine množic in osnovne operacije unija, presek, vsota in produkt na družinah.

\section{Množice in razredi}
\label{sec:mnozice-in-razredi}

V razdelku~\ref{sec:podmnozice} smo spoznali podmnožice, ki jih tvorimo z zapisom
%
\begin{equation*}
    \set{ x \in A \such \phi(x) }.
\end{equation*}
%
Tako podmnožico lahko tvorimo le, če smo predhodno že tvorili njeno nadmnožico~$A$.
%
Mar ne bi bilo bolje dovoliti kar zapis
%
\begin{equation*}
  \set{x \such \phi(x) },
\end{equation*}
%
s katerim bi tvorili množico vseh elementov, ki zadoščajo pogoju~$\phi$?
%
Čemu služi omejitev na elemente~$A$?

Na to vprašanje je odgovoril Bertrand Russell.\footnote{Bertrand Arthur William Russell (1872--1970), znameniti filozof, logik in matematik, prejemnik Nobelove nagrade za književnost.}
%
Razmislek, ki se po njem imenuje \df{Russellov paradoks}, je v matematiko vnesel pravo krizo temeljev, iz katere je v prvi polovici 20.~stoletja izšla logika in temelji matematike, kot jih poznamo danes.
%
Njegov paradoks gre takole. Denimo, da bi lahko tvorili množico vseh množic, ki niso element same
sebe:
%
\begin{equation*}
    R \defeq \set{ S \such \text{$S$ je množica in $S \not\in S$} }.
\end{equation*}
%
Tedaj bi lahko izpeljali protislovje, ker bi hkrati veljalo $R \not\in R$ in $R \in R$:
%
\begin{itemize}
\item dokažimo $R \not\in R$:
  %
  če bi veljalo $R \in R$, bi po definiciji $R$ veljalo $R \not\in R$, kar
  je v protislovju s predpostavko $R \in R$;

\item
  dokažimo $R \in R$: dokazali smo že $R \not\in R$, torej po definiciji $R$ velja $R \in R$.
\end{itemize}
%
Nekaj je narobe z zamislijo, da bi tvorili množice naivno z zapisom $\set{x \such \phi(x)}$.

V sodobni teoriji množic Russellov paradoks razrešimo tako, da ločimo med dvema
zvrstema zbirk ali skupkov elementov, namreč \df{množicami} in \df{razredi}:
%
\begin{enumerate}
\item \df{urelementi} so matematični objekti, ki niso skupki,
\item \df{množice} so skupki, katerih elementi so urelementi in množice,
\item \df{razredi} so skupki, katerih elementi so urelementi in množice.
\end{enumerate}
%
Na prvi pogled ni razlike med množicami in razredi, a to ni res:
\emph{množica je lahko element množice ali razreda,
razred pa ne more biti element množice ali razreda.}
%
Zapis
%
\begin{equation*}
    x \in Y
\end{equation*}
%
je veljaven v primeru, da je $x$ urelement ali množica in $Y$ množica ali razred.
Če je $x$ razred, tedaj $x \in Y$ ni veljavna izjava, podobno kot $1/0$ ni veljaven izraz, in sploh ne moremo govoriti o njegovi resničnostni vrednosti.
%
Povejmo pravilo, s katerim tvorimo razrede.

\begin{pravilo}[Razredi]
  \label{pravilo:tvorba-razredov}
  Za vsako formulo $\phi(x)$ je $\set{ x \such \phi(x) }$ razred, katerega elementi
  so vsi urelementi in množice, ki zadoščajo pogoju~$\phi$:
  %
  \begin{equation*}
    a \in \set{ x \such \phi(x) } \liff \phi(a).
  \end{equation*}
\end{pravilo}

\noindent
%
Predpišimo tudi pravilo ekstenzionalnosti za razrede.

\begin{pravilo}[Ekstenzionalnost razredov]
  \label{pravilo:ekstenzionalnost-razredov}
  Razreda~$C$ in~$D$ sta enaka, ko imata iste elemente:
  %
  $
    C = D \liff \all{x} (x \in C \liff x \in D)
  $.
\end{pravilo}

\begin{zgled}
  \df{Russellov razred} $R \defeq \set{ S \such S \not\in S }$ vsebuje vse množice, ki niso element same sebe. Paradoks smo razrešili, saj je nesmiselno zapisati $R \in R$.
\end{zgled}

\begin{zgled}
  \df{Razred vseh množic}
  %
  \begin{equation*}
    \Set \defeq \set{ S \such \text{$S$ je množica} },
  \end{equation*}
  %
  ki ga nekateri avtorji označijo tudi z $V$, vsebuje vse množice.
\end{zgled}

\begin{zgled}
  Tvorimo lahko tudi razrede matematičnih struktur, na primer razred vseh vektorskih prostorov $\set{V \such \text{$V$ je vektorski prostor}}$, razred vseh grup itn.
\end{zgled}

\begin{zgled}
  Naj bo $S$ množica. Razred $\set{x \such x \in S}$ vsebuje natanko elemente množice~$S$.
\end{zgled}

Zadnji zgled pove, da so nekateri razredi pravzaprav množice.

\begin{definicija}
  Razred $C$ je \textbf{enak množici}~$S$, če imata iste elemente, $\all{x} x \in C \liff x \in S$, v tem
  primeru pišemo $C = S$. Razred je \textbf{nepravi}, če je enak kaki množici in \textbf{pravi} sicer.
\end{definicija}

\begin{zgled}
  Russellov razred
  $R \defeq \set{ S \such \text{$S$ je množica in $S \not\in S$} }$
  je pravi. Če bi bil množica, bi sledil Russellov paradoks.
\end{zgled}

\begin{zgled}
  Če je pravi razred~$C = \set{ x \such \phi(x)} \subesteq D$ podrazred razreda~$D$, potem je tudi~$D$ pravi razred. Če bi namreč~$D$ bil množica, bi bil tudi~$C$ množica, saj bi ga tvorili kot podmnožico $\set{x \in D \such \phi(x)}$.
\end{zgled}

\begin{zgled}
  Razred vseh množic~$\Set$ je pravi, saj vsebuje pravi podrazred~$R$.
\end{zgled}

Z razredi lahko delamo tako kot z množicami: tvorimo unije, preseke in produkte
razredov in govorimo o podrazredih. Pri tem uporabljamo enake oznake za
operacije kot pri množicah. Paziti moramo le, da razreda nikoli ne uporabimo kot
element kake množice ali razreda. Na primer, če je $C$ razred, lahko tvorimo
potenčni razred $\pow{C}$, ki vsebuje vse \emph{podmnožice} $C$:
%
\begin{equation*}
  \pow{C} \defeq \set{ S \such S \in \Set \land S \subseteq C }.
\end{equation*}

\section{Družine množic}
\label{sec:druzine-mnozic}


Pogosto imamo opravka z zbirko množic, ki je odvisna od parametra. Na primer, množico vseh deliteljev števila $n \in \NN$ definiramo takole:
%
\begin{equation*}
  D_n \defeq \set{m \in \NN \such \some{k \in \NN} n = k \cdot m}.
\end{equation*}
%
V zgornjem zapisu lahko~$D$ obravnavamo kot preslikavo iz naravnih števil v množice.

\begin{definicija}
  \label{def:druzina-mnozic}
  \df{Družina množic} je preslikava $I \to \Set$. Množici $I$ pravimo \df{indeksna množica} in njenim elementom \df{indeksi}.
\end{definicija}


\begin{zgled}
  Zgornji primer je družina množica $D : \NN \to \Set$ z indeksno množico $\NN$ in preslikavo
  %
  \begin{equation*}
    D : n \mapsto \set{m \in \NN \such \some{k \in \NN} n = k \cdot m}.
  \end{equation*}
\end{zgled}

\begin{zgled}
  Končno zbirko množic lahko indeksiramo s končno množico. Denimo, da imamo
  množice $\set{2, 4, 6}$, $\set{1, 3, 5}$, $\set{10}$, $\set{7, 8, 9}$ in $\set{0}$. Iz njih lahko tvorimo družino $S : \set{1, 2, 3, 4, 5} \to \Set$:
  %
  \begin{align*}
  S_1 &\defeq \set{2, 4, 6},\\
  S_2 &\defeq \set{1, 3, 5},\\
  S_3 &\defeq \set{10},\\
  S_4 &\defeq \set{7, 8, 9},\\
  S_5 &\defeq \set{0}.
  \end{align*}
\end{zgled}

\begin{zgled}
  Nihče nas ne sili, da morajo biti indeksi števila. V prejšnjem primeru bi lahko
  definirali $S : \set{42, 13, \sqrt{2}, \emptyset, \RR} \to \Set$
  %
  \begin{align*}
  S_{42} &\defeq \set{2, 4, 6},\\
  S_{13} &\defeq \set{1, 3, 5},\\
  S_{\sqrt{2}} &\defeq \set{10},\\
  S_{\emptyset} &\defeq \set{7, 8, 9},\\
  S_{\RR} &\defeq \set{0}.
  \end{align*}
\end{zgled}

\begin{zgled}
  Zaprti intervali tvorijo družino $S : I \to \Set$, kjer je
  %
  \begin{align*}
    I &\defeq \set{(a, b) \in \RR \times \RR \such a < b}, \\
    S_{(a, b)} &\defeq \set{x \in \RR \such a \leq x \leq b}.
  \end{align*}
\end{zgled}

\begin{zgled}
  \label{zgled:druzina-premic}
  Premico v ravnini, ki gre skozi izhodišče in z osjo~$x$ oklepa kot~$\alpha$, opišemo z množico
  %
  \begin{equation*}
    P_\alpha \defeq \set{(x, y) \in \RR \times \RR \such x \sin \alpha = y \cos \alpha}.
  \end{equation*}
  %
  Tako dobimo družino $P : \RR \to \Set$, katere indeksi so koti.
\end{zgled}

\begin{zgled}
  Množice v družini se lahko ponavljajo. V družini premic iz zgleda~\ref{zgled:druzina-premic} velja na primer $P_{\pi/4} = P_{5\pi/4}$ in $P_0 = P_\pi$. V splošnem za vsak $\alpha \in \RR$ in $k \in \ZZ$ velja $P_{\alpha + k \pi} = P_\alpha$, zato se vsaka premica pojavi neskončnokrat. Skrajni primer je \df{konstantna družina}, v kateri so vse množice med seboj enake.
\end{zgled}

\begin{zgled}
  Besedi \nls{prazna} in \nls{neprazna} se pri družinah lahko nanašata na indeksno množico ali na množice v družini.
  %
  \df{Prazna družina} je družina množic, ki je indeksirana z $\emptyset$.
  %
  Ločiti jo moramo od \df{družine praznih množic} $I \to \Set$, $i \mapsto \emptyset$.
  %
  Podobno je \df{neprazna družina} indeksirana z neprazno množico, v \df{družini nepraznih množic} pa so vse množice neprazne.
\end{zgled}


\section{Konstrukcije in operacije z družinami množic}
\label{sec:operacije-z-druzinami}

Operacije $\times$, $+$, $\cap$ in $\cup$ lahko posplošimo tako, da namesto z dvema
množicama delujejo na poljubnem številu množic. V ta namen uporabimo družine
množic.

\subsection{Presek in unija družine}
\label{sec:presek-in-unija-druzine}

Presek in unija družine $A : I \to \Set$ sta definirana takole:
%
\begin{align*}
  \Union{i \in I} A_i &\defeq \set{ x \such \some{i \in I} x \in A_i },
  \\
  \Inter{i \in I} A_i &\defeq \set{ x \such \all{ i \in I} x \in A_i }.
\end{align*}
%
Pozor! Na desni strani imamo razred! Res se lahko zgodi, da dobimo pravi razred, namreč
kot presek prazne družine:
%
\begin{align*}
  \Inter{i \in \emptyset} A_i
  &= \set{ x \such \all{ i \in \emptyset} x \in A_i } \\
  &= \set { x \such \top } \\
  &= \Set.
\end{align*}
%
Kdaj pa dobimo množico? Presek neprazne družine je vedno množica. Res, če imamo
$k \in I$, potem velja
%
\begin{equation*}
    \Inter{i \in I} A_i = \set{ x \in A_k \such \all{i \in I} x \in A_i}.
\end{equation*}
%
Sedaj na desni ne stoji več razred, ampak podmnožica množice~$A_k$.

Kaj pa unija družine množic? Ali je množica? Izkaže se, da za to potrebujemo dodatno pravilo.

\begin{pravilo}[Unija družine]
  \label{pravilo:unija-druzine}
  Unija družine množic je množica.
\end{pravilo}


\subsection{Kartezični produkt družine}
\label{sec:produkt-druzine}

\begin{definicija}
  \label{def:funkcija-izbire}
  \df{Funkcija izbire} za družino $A : I \to \Set$ je tako prirejanje, ki vsakemu indeksu $i \in I$ priredi natanko en element $f(i) \in A_i$.
\end{definicija}

\begin{zgled}
  Primer: funkcija izbire za družino
  %
  \begin{align*}
  A &: \NN \to \Set \\
  A_n &\defeq \set{ x \in \RR \such 0 < x < 2^{-n} }
  \end{align*}
  %
  je na primer $f(n) \defeq 2^{-n - 1}$. To ni edina funkcija izbire za $A$, lahko bi vzeli tudi $f(n) \defeq 2^{-n} / 3$.
\end{zgled}

\begin{definicija}
  \label{def:produkt-druzine}
  \df{Kartezični produkt} družine $A : I \to \Set$ je množica vseh funkcij izbire družine~$A$:
  %
  \begin{equation*}
    \Prod{i \in I} A_i \defeq
    \set{f : I \to \Union{i \in I} A_i \such \all{i \in I} f(i) \in A_i}.
  \end{equation*}
  %
  Za vsak $j \in I$ imamo \df{$j$-to projekcijo}
  %
  \begin{align*}
    \pr[j] &:  \left(\Prod{i \in I} A_i\right) \to A_j, \\
    \pr[j] &:  f \mapsto f(j).
  \end{align*}
\end{definicija}

Običajni kartezični produkt dveh množic je poseben primer produkta množic, namreč družine
množic, ki je indeksirana z $I = \set{1, 2}$. Natančneje, velja
%
\begin{equation*}
  A \times B \cong \Prod{i \in \set{1, 2}} C_i,
\end{equation*}
%
kjer je $C_1 \defeq A$ in $C_2 \defeq B$.

Tudi eksponentna množica je poseben primer produkta množic, saj velja
%
\begin{equation*}
  B^A \cong \Prod{a \in A} B.
\end{equation*}
%
Na desni imamo produkt konstantne družine množic
%
\begin{align*}
  A &\to \Set, \\
  a &\mapsto B.
\end{align*}


\subsection{Koprodukt ali vsota množic}
\label{sec:koprodukt-druzine}

Vsoto množic posplošimo na koprodukt družine.

\begin{definicija}
  \label{def:koprodukt-druzine}
  \df{Koprodukt} ali \df{vsota družine} $A : I \to \Set$ je množica
  $\Sum{i \in I} A_i$, katere elementi so $\inj[i](a)$ za $i \in I$ in $a \in A_i$.
  Preslikavi $\inj[k] : A_k \to \Sum{i \in I} A_i$ pravimo \df{$k$-ta injekcija}.

  Poleg tega definiramo še \df{projekciji}
  %
  \begin{align*}
    \fst (\inj[i](a)) &\defeq i, \\
    \snd (\inj[i](a)) &\defeq a.
  \end{align*}
  %
  Namesto $\sum$ se piše tudi $\coprod$.
\end{definicija}

Poseben primer koprodukta je vsota $A + B$, saj velja
%
\begin{equation*}
  A + B \cong \Sum{k \in \set{1, 2}} C_k,
\end{equation*}
%
kjer je
%
\begin{align*}
  C &: \set{1, 2} \to \Set \\
  C_1 &\defeq A, \\
  C_2 &\defeq B.
\end{align*}
%
Tudi kartezični produkt $A \times B$ je poseben primer koprodukta, saj velja
%
\begin{equation*}
    A \times B \cong \Sum{a \in A} B.
\end{equation*}
%
Na desni imamo tokrat koprodukt konstantne družine množic
%
\begin{align*}
  A &\to \Set, \\
  a &\mapsto B.
\end{align*}
